\documentclass[11pt,a4paper]{article}

% --- Essential Mathematical & Physical Packages ---
\usepackage[utf8]{inputenc}
\usepackage[T1]{fontenc}
\usepackage{amsmath,amssymb,amsfonts,amsthm,mathtools}
\usepackage{bm}
\usepackage{physics}
\usepackage{geometry}
\geometry{top=25mm,bottom=25mm,left=25mm,right=25mm}

% --- Typography & Layout ---
\usepackage{microtype}
\usepackage{booktabs}
\usepackage{array}
\usepackage{enumitem}
\usepackage{xcolor}
\usepackage{hyperref}
\hypersetup{
    colorlinks=true,
    linkcolor=blue!70!black,
    citecolor=red!70!black,
    urlcolor=blue!60!black
}

% --- Theorem & Axiom Environments ---
\newtheorem{axiom}{Axiom}
\newtheorem{theorem}{Theorem}
\newtheorem{lemma}{Lemma}
\newtheorem{proposition}{Proposition}
\theoremstyle{definition}
\newtheorem{definition}{Definition}
\theoremstyle{remark}
\newtheorem{remark}{Remark}

% --- Standard Mathematical Operators ---
\newcommand{\Tr}{\mathrm{Tr}}
\newcommand{\Area}{\mathrm{Area}}
\newcommand{\supp}{\mathrm{supp}}
\newcommand{\RealPart}{\mathrm{Re}}
\newcommand{\ImagPart}{\mathrm{Im}}

\title{\textbf{A Continuum-Mechanical and Non-Equilibrium Thermodynamic Framework of Physical and Biological Existence}}
\author{\textbf{Ishan Tomar} \\ \textit{Vidyaman Research Institute}}
\date{\today}

\begin{document}

\maketitle

\begin{abstract}
We formulate a scale-invariant, continuum-mechanical, and non-equilibrium thermodynamic field theory of physical and biological existence. Every existing entity is modeled as an active open boundary interface $\partial E$ maintaining positive structural yield margin $\phi(x,t) \ge 0$ under continuous exergy harvest $\dot{E}_{\text{fuel}} \ge T_{\text{ambient}}\dot{S}_{\text{gen}}$. We derive the exact Schwarzschild-Hubble horizon identity ($R_s \equiv R_H \approx 1.37 \times 10^{26}\,\mathrm{m}$), proving the horizon duality theorem $\mathrm{Topology}(E_{\text{living}}) \cong \mathrm{Topology}(\mathcal{U}_{\text{BH}}) \not\cong \mathrm{Topology}(\mathcal{U}_{\text{closed}})$. Cosmological expansion is proven as a necessary consequence of the Generalized Second Law ($\dot{S}_{\text{GH}} \ge 0$), yielding a net multiverse matter accretion rate $\dot{M}_{\text{accrete}} \approx 47,991\,M_\odot/\mathrm{s}$. Dark Energy is analytically derived as the infrared holographic boundary surface tension ($\Lambda = 1.091 \times 10^{-52}\,\mathrm{m^{-2}}$, matching observations to $1.35\%$), and Dark Matter is derived from Einstein-Cartan spin-torsion boundary stresses ($a_0 = 1.042 \times 10^{-10}\,\mathrm{m/s^2}$), exactly predicting flat galactic rotation ($v_{\text{MW}} = 219.7\,\mathrm{km/s}$ to $0.14\%$). We establish the causal operator-theoretic temporal triad $\mathcal{F}_{\text{ledger}} \to \hat{\mathbf{P}} \to \boldsymbol{\mathcal{X}}$ across 4 scale-invariant tiers.
\end{abstract}

\vspace{1em}
\tableofcontents
\newpage


\section*{Cognitive Thermodynamics, The Ariṣaḍvarga, and the Landauer Veto}

\begin{quote}\textbf{Cognitive \& Psychological Focus:} This treatise formalizes \textbf{Tier III Cognitive Architectures}, \textbf{Bayesian Active Inference Engines}, the \textbf{Thermodynamics of the Six Internal Afflictions (\textit{Ariṣaḍvarga})}, and the \textbf{Landauer Information-Erasure Cost of Voluntary Veto (\textit{Viveka})}.\end{quote}

---

\section{1. Tier III: Cognitive Agents as Complex Predictive Engines}

In [`draft.md` §2.3.4 \& §4](file:///c:/Users/tomar/Documents/Vidyaman/Project_writeup_1/essays/existence/draft.md), a cognitive agent ($E \in T_{\text{III}}$) operates across complexified phase space $\Omega_{\mathbb{C}} = \Omega_{\mathbb{R}} \oplus i \Omega_{\mathfrak{Im}}$:
* \textbf{Real Substrate ($\Omega_{\mathbb{R}}$):} Immediate somatic, motor, and neurochemical configurations.
* \textbf{Imaginary Substrate ($\Omega_{\mathfrak{Im}}$):} Internal generative Bayesian models simulating counterfactual future trajectories ($\hat{\mathbf{C}}_{\text{future}}$).

\begin{verbatim}
                         THE COGNITIVE ENERGY ALLOCATION FORK
                                          │
                         Total Informational Fuel: Ė_total
                                          │
             ┌────────────────────────────┴────────────────────────────┐
             │                                                         │
[REAL MOTOR ACTUATION: Ė_Re]                               [IMAGINARY SIMULATION: Ė_Im]
• Immediate somatic response                               • Counterfactual future forward simulation
• Physical muscular work                                   • Sagawa-Ueda negentropy bound:
• χ* ≈ 1: Optimal balance                                  • ⟨W_dissipated⟩ ≥ Δℱ_noneq - k_B T Δ𝓘(Ĉ; C)
\end{verbatim}

\subsection{The Sagawa-Ueda Information-Thermodynamic Bound:}
The dissipation required to survive an environmental shock is reduced by the mutual information $\Delta \mathcal{I}(\hat{\mathbf{C}}; \mathbf{C}_{\text{future}})$ acquired via internal simulation:

\begin{equation}
\boxed{\langle W_{\text{dissipated}} \rangle \ge \Delta \mathcal{F}_{\text{noneq}} - k_B T \cdot \Delta \mathcal{I}\left( \hat{\mathbf{C}}_{\text{predicted}}; \, \mathbf{C}_{\text{future}} \right)}
\end{equation}

Pre-stiffening cognitive priors ($\mathbf{R}_{\text{active}}(t - \Delta t)$) reduces physical damage to near zero upon predictable impacts.

---

\section{2. Involuntary Latency (\textit{Vāsanā}) vs. Voluntary Landauer Veto Gating (\textit{Viveka})}

Living cognitive agents possess two distinct modes of interfacial actuation:

\begin{verbatim}
                            THE ACTUATION PATHWAYS
                                       │
        ┌──────────────────────────────┴──────────────────────────────┐
        │                                                             │
[INVOLUNTARY REFLEX: VĀSANĀ]                     [VOLUNTARY VETO: VIVEKA]
• Fast feed-forward reactive arc                 • Slow Dyson time-ordered loop
• Response latency: τ ~ 10^(-2) s                • Pays Landauer erasure cost:
• Computation entropy: Ḣ ≈ 0                    • Ė_veto = k_B T ln 2 · Ḣ_erasure > 0
• Automated habit execution                      • Active inhibitory gating 𝓞_inhibit
\end{verbatim}

1. \textbf{Involuntary Reflex (\textit{Vāsanā / Saṃskāra}):}  
   An automated feed-forward reflex arc where incoming challenge traction $\mathbf{C}$ directly triggers motor expression $\boldsymbol{\mathcal{X}}$ through pre-configured neural pathways ($\tau_{\text{latency}} \sim 10^{-2}\,\mathrm{s}$). Because no deliberative computation or state re-evaluation occurs, the Landauer informational entropy cost is zero ($\dot{\mathcal{H}} \approx 0$).

2. \textbf{Voluntary Discriminative Veto (\textit{Viveka}):}  
   The conscious intervention of the higher-order Dyson loop ($\mathcal{T}\exp(\int \hat{\mathcal{L}} d\tau)$), which inspects the automated trajectory and applies an active inhibitory gate ($\mathcal{O}_{\text{inhibit}}$) to prevent expression.
   * \textbf{The Landauer Metabolic Tax:} Erasing the automated motor command requires physical bit erasure, demanding continuous free-energy dissipation:
     
\begin{equation}
\boxed{\dot{\mathcal{E}}_{\text{veto}} = k_B T \ln 2 \cdot \dot{\mathcal{H}}_{\text{erasure}} > 0}
\end{equation}

   \textit{ Exercising conscious restraint (}Viveka / Dama*) is physically exhausting because it consumes real biochemical glucose ($k_B T \ln 2$) to clear neural register buffers.

---

\section{3. Top-Down Craving (\textit{Kāma}) vs. Bottom-Up Hunger (\textit{Kṣudhā})}

\textit{ \textbf{Hunger (}Kṣudhā*):}  
  A bottom-up First-Law exergy deficit occurring entirely in real configuration space $\Omega_{\mathbb{R}}$:
  
\begin{equation}
\dot{E}_{\text{fuel}} < \dot{E}_{\text{crit}} \equiv T_{\text{ambient}} \int_E \sigma_{\text{total}} \, dV \implies \frac{d\mathcal{G}}{dt} < 0
\end{equation}

  \textit{Kṣudhā} is a biochemical indicator of metabolic depletion (falling blood glucose, low ATP/AMP ratio) requiring zero memory simulation.

\textit{ \textbf{Desire / Craving (}Kāma / Tṛṣṇā*):}  
  A top-down simulation executed in imaginary phase space $\Omega_{\mathfrak{Im}}$:
  
\begin{equation}
D_{\mathfrak{Im}} = \hat{\pi}(\mathcal{F}_{\text{ledger}}) \xrightarrow{\text{simulation}} \Delta \hat{\mathcal{G}}_{\text{predicted}} > 0
\end{equation}

  \textit{Kāma} occurs when the cognitive model recalls past reward states from $\mathcal{F}_{\text{ledger}}$ and projects an artificial free-energy surplus. This top-down prediction transduces downward into physical physiology via the Semantic Transduction Tensor ($\mathbf{C}_{\text{craving}} = \mathbf{K}_{\text{trans}} \cdot \nabla_\theta \mathcal{D}_{\text{KL}}$), generating visceral dopamine cascades and restlessness \textbf{even when fully satiated}.

---

\section{4. Thermodynamics of the \textit{Ariṣaḍvarga} (The Six Internal Afflictions)}

In classical Sanskrit psychology, the \textit{Ariṣaḍvarga} represents the six inner enemies that destroy discernment. In our framework, these are derived as \textbf{exact mathematical distortions of the structural margin tensor $\phi(x, t)$}:

\begin{verbatim}
                       THE ARIṢAḌVARGA TENSOR DISTORTIONS
                                       │
     ┌─────────────────────────────────┼─────────────────────────────────┐
     │                                 │                                 │
[KĀMA: PHANTASMAGORIC OVERDRIVE]  [KRODHA: OVER-REACTIVE STRESS]    [LOBHA: RESOURCE STRANGULATION]
Ė_Im ⟶ ∞, Ė_Re ⟶ 0                ||R|| >> ||C|| (Self-Fracture)   J_metabolic ⟶ 0 (Local Gluttony)
     │                                 │                                 │
[MOHA: INVERTED PRIORS]           [MADA: HUBRISTIC BLINDNESS]       [MĀTSARYA: MALICIOUS SABOTAGE]
||Ĉ - C|| ⟶ ∞ (Boundary Delusion) Ĉ ≡ 0 (Brittle Impact Failure)   Consuming fuel to drop Δϕ_AB
\end{verbatim}

1. \textbf{Kāma (काम — Unbounded Imaginary Projection Overdrive):}  
   Informational fuel allocation diverges ($\chi \to \infty \implies \dot{\mathcal{E}}_{\mathfrak{Im}} \to \dot{\mathcal{E}}_{\text{total}}, \dot{\mathcal{E}}_{\mathfrak{Re}} \to 0$). The entity expends all free energy simulating ungrounded imaginary prospects while the physical substrate starves into anergic boundary collapse.
2. \textbf{Krodha (क्रोध — Explosive Over-Reactive Stress Discharge):}  
   Internal counter-stress vastly exceeds external challenge ($\|\mathbf{R}_{\text{active}}\| \gg \|\mathbf{C}\|$). The localized stress concentration generates internal Rankine-Hugoniot shock waves that fracture the entity's own boundary from within.
3. \textbf{Lobha (लोभ — Parasitic Resource Hoarding):}  
   A constituent node hoards fuel measure ($\mu(E^j) \uparrow$) while choking metabolic syncytial flux ($\mathbf{J}_{\text{metabolic}}^{j \to k} \to 0$). This starves neighboring nodes, triggering a localized ischemia cascade that collapses the macro-envelope.
4. \textbf{Moha (मोह — Perceptual Delusion / Model Inversion):}  
   The generative model divergence explodes ($\|\hat{\mathbf{C}} - \mathbf{C}_{\text{true}}\| \to \infty$). The entity mistakes destructive threats for fuel sources, actively opening its boundary pores ($\mathbf{v}_n \cdot \hat{n} < 0$) to toxic foreign challenge fields.
5. \textbf{Mada (मद — Hubristic Invulnerability / Zero-Challenge Prior):}  
   The entity sets its predicted external challenge to zero ($\hat{\mathbf{C}} \equiv \mathbf{0}$), disarming its adaptive viscoelastic dampers ($\nu \to 0$). Upon the arrival of real environmental shocks, it undergoes catastrophic brittle fracture without damping.
6. \textbf{Mātsarya (मात्सर्य — Malicious Zero-Sum Trophic Sabotage):}  
   An entity expends its own metabolic exergy solely to diminish a neighbor's margin ($\Delta \phi_{AB} < 0$), lowering total syncytial free energy and precipitating mutual systemic lysis.


\end{document}
